\par
\SourceLeaf{378}%
Wir glauben, noch einige Bemerkungen hinzufügen zu müssen. Bei sämtlichen Messungen wurde der Index zu Beginn jeder Reihe stets auf 0 zurückgestellt. Dieses Verfahren können wir nicht gutheißen, weil dadurch, so sehr man auch die Zahl der Reihen erhöht, stets derselbe vom Teilungsfehler abhängige konstante Fehler im Ergebnis verbleiben muss. --- Bei den Messungen der piemontesischen Offiziere ist jedes Mal der Zustand der Luft angegeben. Unter 414 Messreihen zählen wir 320, bei denen Windstille, und 94, bei denen Wind vermerkt ist. Ein so günstiges Verhältnis hätte man an so hoch gelegenen Standpunkten kaum erwartet (der höchste liegt 3534 Meter über dem Meeresspiegel).

Der zweite Band enthält in zehn Abschnitten die Arbeiten der Astronomen. Die ersten beiden Abschnitte enthalten die Bestimmungen von Längenunterschieden durch Pulversignale für die Längengradmessung. Die ersten Versuche dieser Art fanden im September 1821 statt. Die Pulversignale wurden auf der Rocca Melone gegeben und an der 170,000 Meter entfernten Mailänder Sternwarte und auf dem nahen Mont Cenis beobachtet. Für die Zeitbestimmung am letzteren Ort war im Garten des Hospizes eine kleine Sternwarte errichtet und darin ein Passagenfernrohr von \textsc{Fortin} aufgestellt worden. Es scheint jedoch nicht von hervorragender Qualität gewesen zu sein, wie für die Zapfen, einen wesentlichen Teil des Instruments, ausdrücklich bemerkt wird. Die Rocca Melone war von hier aus nicht sichtbar. Um die Signale zu sehen, musste man eine etwas entfernte Stelle aufsuchen und die Zeit mit einem Chronometer von \textsc{Earnshaw} dorthin übertragen. Auch die Beobachtungen am Passagenfernrohr wurden meist nach diesem Chronometer notiert. Die Aufzeichnungen machte ein Gehilfe auf ein Zeichen des Beobachters hin. Alle diese Umstände mindern freilich zusammengenommen das Vertrauen in die. Genauigkeit des Endergebnisses, auch wenn die drei Teilergebnisse der drei Beobachtungstage sehr gut übereinstimmen. Dazu kommt, dass man hier die Zeit aus Sternen bestimmte, in Mailand dagegen aus Sonnendurchgängen, und schließlich, dass man beim Passagenfernrohr von \textsc{Fortin} den rechten Winkel zwischen der optischen Achse und der Drehachse offenbar gar nicht justierte. Zumindest wird dieser wichtige Umstand bei diesen Beobachtungen überhaupt nicht erwähnt.

Bei den ähnlichen Arbeiten im Jahr 1822 ging man
